\chapter{Implicit Differentiation}

Most curves in the plane are not graphs. The circle $x^{2}+y^{2}=25$ fails
the vertical line test, the equation $\mathrm{e}^{y}+xy=1$ cannot be solved
for $y$ in closed form at all, and an indifference curve is handed to us as
a level set rather than as a formula. Yet each of them has a tangent at
almost every point, and economics asks for its slope constantly: that slope
\emph{is} the marginal rate of substitution.

Implicit differentiation gets it without solving for $y$. The move rests on
one idea, and the whole chapter is that idea applied repeatedly: pretend the
equation has already been solved near the point of interest, so that $y$ is
a function $y(x)$; substituting it makes the equation an identity in $x$;
identities may be differentiated. Because every $y$ in sight is secretly the
composite $y(x)$, the chain rule leaves a factor $y'$ behind wherever a $y$
was, and the resulting linear equation is solved for $y'$.

Two questions have to be kept apart, and the chapter answers both. \emph{Is}
there such a function $y(x)$ near the point --- and if so, what is its
slope? The first question is the Implicit Function Theorem, and the honest
answer is sometimes no: at a point where the curve crosses itself or stands
vertically, no function $y(x)$ describes it, and the slope formula returns a
division by zero to say so.

\section{Core Rules}

\begin{rulebox}{Why a factor $y'$ appears}
Suppose that near a point of the curve the equation really does pin down one
$y$ for each $x$, and call that function $y(x)$. Substituting it turns the
equation into an \textbf{identity in $x$}, true for every $x$ near $x_{0}$:
\[
  x^{2}+\bigl(y(x)\bigr)^{2}=25 .
\]
An identity may be differentiated on both sides. Every $y$ in the expression
is the composite $y(x)$, so it is differentiated by the \textbf{chain rule}
and leaves a factor $y'$ behind:
\[
  \frac{\mathrm{d}}{\mathrm{d}x}\bigl[y^{2}\bigr]
  =\frac{\mathrm{d}}{\mathrm{d}x}\bigl[(y(x))^{2}\bigr]
  =2\,y(x)\,y'(x)=2yy' .
\]
Hence $2x+2yy'=0$ and $y'=-x/y$. Nothing was solved for $y$, and the answer
is \emph{allowed} to contain $y$: it is the slope at a point of the curve,
and the two points with the same $x$ --- here $(3,4)$ and $(3,-4)$ --- have
slopes $-3/4$ and $+3/4$.
\end{rulebox}

\begin{rulebox}{Setup}
Suppose $F(x,y)=0$ defines $y=y(x)$ implicitly. Differentiate every term
with respect to $x$. Whenever an expression involves $y$, the chain rule
produces an extra factor $\dd{y}{x}$:
\[
  \frac{\mathrm{d}}{\mathrm{d}x}\bigl[g(y)\bigr]=g'(y)\cdot\dd{y}{x} .
\]
So $\frac{\mathrm{d}}{\mathrm{d}x}[\sin y]=\cos(y)\,y'$,
$\frac{\mathrm{d}}{\mathrm{d}x}[\mathrm{e}^{y}]=\mathrm{e}^{y}y'$, and
$\frac{\mathrm{d}}{\mathrm{d}x}[\ln y]=y'/y$.
\end{rulebox}

\begin{rulebox}{Mixed terms}
For a product such as $x^{3}y^{2}$ apply the product rule, treating $y$ as
a function of $x$:
\[
  \frac{\mathrm{d}}{\mathrm{d}x}\bigl[x^{3}y^{2}\bigr]=3x^{2}y^{2}+x^{3}\cdot 2yy' .
\]
\end{rulebox}

\begin{rulebox}{Solving for $y'$}
Collect all terms containing $y'$ on one side and factor:
\[
  A(x,y)\,y'+B(x,y)=0\;\Longrightarrow\;y'=-\frac{B(x,y)}{A(x,y)} .
\]
The equation is always \emph{linear} in $y'$, however complicated it looks,
because $y'$ only ever enters through one chain-rule step.
\end{rulebox}

\begin{rulebox}{The curve is a level set --- and that gives the formula at once}
Write the equation as $F(x,y)=0$. Then the curve is the \textbf{level set of
the surface $z=F(x,y)$ at height $0$}: the shoreline where the surface cuts
the horizontal plane $z=0$. Walking along the curve, $F$ never changes.
Writing that down with the chain rule for $F\bigl(x,y(x)\bigr)=0$,
\[
  \pd{F}{x}+\pd{F}{y}\cdot\dd{y}{x}=0
  \qquad\Longrightarrow\qquad
  \boxed{\;\dd{y}{x}=-\frac{F_{x}(x,y)}{F_{y}(x,y)}\;}
  \qquad\text{provided } F_{y}(x,y)\neq0 .
\]
This is not a new method, it is the same computation done once and for all:
for $F=x^{2}+y^{2}-25$ it returns $-2x/2y=-x/y$, as above. In economics the
curve is an indifference curve $U(x,y)=\bar u$ and the formula reads
$\dd{y}{x}=-U_{x}/U_{y}$: the marginal rate of substitution is the ratio of
marginal utilities.
\end{rulebox}

\begin{gallery}
\begin{tikzpicture}
  \begin{axis}[surfaxis, title={$z=F(x,y)=x^{3}-3x^{2}+y^{2}-2y+1$}]
    \addplot3[surf, domain=-1.2:3.2, y domain=-1.6:3.6]{x^3-3*x^2+y^2-2*y+1};
  \end{axis}
\end{tikzpicture}\hfill
\begin{tikzpicture}
  \begin{axis}[geoaxis, xmin=-1.6, xmax=3.6, ymin=-1.9, ymax=3.9,
      xtick={-1,1,2,3}, ytick={-1,1,2,3}, title={the level set $F=0$}]
    \addplot[domain=-1.2:3, samples=241]{1+abs(x)*sqrt(3-x)};
    \addplot[domain=-1.2:3, samples=241]{1-abs(x)*sqrt(3-x)};
    \addplot[only marks, mark=*, mark size=1.3pt, accent, forget plot]
      coordinates {(0,1) (3,1)};
  \end{axis}
\end{tikzpicture}
\gcap{The surface and its zero level set $k$. At $(0,1)$ the curve crosses
itself; at $(3,1)$ it stands vertical. Everywhere else a local $y(x)$
exists.}
\end{gallery}

\begin{rulebox}{When there is no $y(x)$: the Implicit Function Theorem}
\textbf{Theorem.} If $F$ is $C^{1}$ near $(x_{0},y_{0})$, if
$F(x_{0},y_{0})=0$, and if
\[
  F_{y}(x_{0},y_{0})\neq0 ,
\]
then there are neighbourhoods on which the equation $F(x,y)=0$ defines
exactly one differentiable function $y=y(x)$ with $y(x_{0})=y_{0}$, and its
derivative is $-F_{x}/F_{y}$.

The condition is exactly what the formula needs, and its failure is
visible in the picture. Where $F_{y}=0$:
\begin{itemize}[leftmargin=1.6em,itemsep=0.3ex]
  \item if $F_{x}\neq0$, the curve has a \textbf{vertical tangent}. No
        function $y(x)$ can describe it there --- but the roles may be
        swapped: $x=x(y)$ exists, with
        $\dd{x}{y}=-F_{y}/F_{x}$, and that derivative is $0$.
  \item if $F_{x}=0$ as well, the point is \textbf{singular}: typically the
        curve crosses itself, and \emph{two} branches with two different
        slopes pass through. No single function describes either the curve
        or its slope, and $-F_{x}/F_{y}$ degenerates to $0/0$. Separate the
        branches first, then differentiate each one.
\end{itemize}
Reading it off a picture: look for vertical tangents and self-intersections.
Those are the only places the theorem can fail. A failed condition is a
warning, not a proof that no $y(x)$ exists: $y^{3}-x=0$ has $F_{y}=0$ at the
origin, yet $y=x^{1/3}$ solves it for every $x$ (with a vertical tangent, so
not differentiably). Chapter~14 collects such cases.
\end{rulebox}

\begin{rulebox}{Logarithmic differentiation}
If $y=f(x)>0$ is a complicated product, quotient or power, take $\ln$
first:
\[
  \ln y=\ln f(x)
  \;\Longrightarrow\;
  \frac{y'}{y}=\frac{\mathrm{d}}{\mathrm{d}x}\ln f(x)
  \;\Longrightarrow\;
  y'=y\cdot\frac{\mathrm{d}}{\mathrm{d}x}\ln f(x) .
\]
Particularly useful for $y=u(x)^{v(x)}$: $\ln y=v(x)\ln u(x)$. This is
implicit differentiation applied to $\ln y-\ln f(x)=0$, which is why the
factor $y'/y$ appears.
\end{rulebox}

\begin{rulebox}{Higher-order implicit}
To find $y''$, differentiate $y'$ a second time with respect to $x$, again
remembering that any $y$ that appears contributes a factor $y'$. Then
substitute the expression already found for $y'$ to eliminate it.
\end{rulebox}

\begin{rulebox}{Tangent line via implicit slope}
The tangent at $(x_{0},y_{0})$ on $F(x,y)=0$ is
$y-y_{0}=y'(x_{0})\cdot(x-x_{0})$, where $y'(x_{0})$ is computed
implicitly. Both coordinates are substituted into $-F_{x}/F_{y}$ before the
line is written down.
\end{rulebox}

\section{Worked Example}

\begin{worked}{The folium $x^{3}+y^{3}=6xy$: find $y'$, the tangent at $(3,3)$, and every point where the tangent is vertical}
\wstep{Step 1. Confirm the point is on the curve.}
$27+27=54$ and $6\cdot3\cdot3=54$. \checkmark\ Do this first: an implicit
slope formula evaluated at a point that is not on the curve returns a
perfectly plausible number that means nothing.

\wstep{Step 2. Differentiate every term with respect to $x$, treating $y$ as $y(x)$.}
Left side: $\frac{\mathrm{d}}{\mathrm{d}x}\bigl[x^{3}+y^{3}\bigr]
=3x^{2}+3y^{2}y'$ --- the chain rule leaves a factor $y'$ behind wherever a
$y$ was.

Right side: $6xy$ is a \textbf{product} of two functions of $x$, so the
product rule applies:
\[
  \frac{\mathrm{d}}{\mathrm{d}x}\bigl[6xy\bigr]=6y+6xy' .
\]
\wtrap{The right side is where this problem is lost. Writing $6y$ alone
(treating $y$ as a constant) or $6xy'$ alone (forgetting the product rule)
are both common. Every occurrence of $y$ needs the chain rule, \emph{and}
every product needs the product rule --- they compound rather than replace
each other.}

\wstep{Step 3. Collect the $y'$ terms; the equation is always linear in $y'$.}
\[
  3x^{2}+3y^{2}y'=6y+6xy'
  \;\Longrightarrow\;
  3y^{2}y'-6xy'=6y-3x^{2}
  \;\Longrightarrow\;
  y'\bigl(3y^{2}-6x\bigr)=6y-3x^{2} .
\]
\[
  \boxed{\;y'=\frac{6y-3x^{2}}{3y^{2}-6x}=\frac{2y-x^{2}}{y^{2}-2x}\;}
\]
The answer contains $y$, and it must: two points of the curve can share an
$x$ and have different slopes, so $x$ alone cannot determine $y'$.

\wstep{Step 4. Evaluate at the point.}
\[
  y'(3,3)=\frac{2\cdot3-9}{9-6}=\frac{-3}{3}=-1 ,
\]
so the tangent is
\[
  y-3=-1(x-3)\quad\Longrightarrow\quad y=6-x .
\]
By the symmetry of the equation in $x$ and $y$, the curve is symmetric about
the line $y=x$, and $(3,3)$ lies on that line --- so a slope of $-1$,
perpendicular to it, is exactly what the geometry predicts. \checkmark

\wstep{Step 5. Cross-check with the level-set formula.}
Writing $F(x,y)=x^{3}+y^{3}-6xy$,
\[
  F_{x}=3x^{2}-6y,\qquad F_{y}=3y^{2}-6x,
  \qquad
  y'=-\frac{F_{x}}{F_{y}}=-\frac{3x^{2}-6y}{3y^{2}-6x}=\frac{6y-3x^{2}}{3y^{2}-6x},
\]
the same expression. The two routes --- differentiate term by term, or
apply $-F_{x}/F_{y}$ --- are the same computation; use whichever you find
faster, but know that they must agree.

\wstep{Step 6. Vertical tangents: where the denominator vanishes.}
$y'$ blows up when $F_{y}=3y^{2}-6x=0$, i.e.\ $x=y^{2}/2$. Substituting into
the curve:
\[
  \frac{y^{6}}{8}+y^{3}=6\cdot\frac{y^{2}}{2}\cdot y=3y^{3}
  \;\Longrightarrow\;
  \frac{y^{6}}{8}=2y^{3}
  \;\Longrightarrow\;
  y^{3}\bigl(y^{3}-16\bigr)=0 .
\]
So $y=0$ (giving $x=0$) or $y=16^{1/3}=2\sqrt[3]{2}\approx2.520$, giving
$x=y^{2}/2\approx3.175$. At the latter the tangent is genuinely vertical.
\wtrap{At the origin the numerator $2y-x^{2}$ also vanishes, so $y'$ is
$0/0$ there --- undefined for a different reason. The folium
\emph{crosses itself} at $(0,0)$, with two distinct tangents, so no single
function $y(x)$ describes the curve near it. That is precisely the case the
Implicit Function Theorem excludes, and a $0/0$ in the slope formula is how
it announces itself.}

\wstep{Step 7. The economic version of the same calculation.}
An indifference curve is a level set $U(x,y)=\bar u$, and the identical
manipulation gives
\[
  \frac{\mathrm{d}y}{\mathrm{d}x}=-\frac{U_{x}}{U_{y}}=-\mathrm{MRS} .
\]
For $U=x^{1/2}y^{1/2}=4$: $U_{x}=\tfrac12x^{-1/2}y^{1/2}$,
$U_{y}=\tfrac12x^{1/2}y^{-1/2}$, so $y'=-y/x$. At $(4,4)$ the slope is
$-1$; at $(8,2)$ it is $-\tfrac14$ --- the indifference curve flattens as
$x$ becomes abundant, which is diminishing marginal rate of substitution
stated as a derivative.
\end{worked}

\section{Practice Problems}

\begin{problems}
  \item Compute $\dd{y}{x}$ implicitly for $x^{2}+4y^{2}=16$.
  \item Find the equation of the tangent line to $x^{2}+4y^{2}=20$ at
        $(2,2)$.
  \item Use logarithmic differentiation to find $y'$ for $y=x^{\sin x}$
        ($x>0$, $\sin x>0$).
  \item Compute $\dd{^{2}y}{x^{2}}$ implicitly for $x^{2}+4y^{2}=4$.
  \item Compute $\dd{y}{x}$ implicitly for $x^{3}-y^{3}=3xy$.
  \item Compute $\dd{y}{x}$ implicitly for $x^{2}y^{3}=8$.
  \item Find the equation of the tangent line to $x^{2}y+y^{2}x=6$ at
        $(1,2)$.
  \item Use logarithmic differentiation to find $y'$ for $y=(\ln x)^{x}$
        ($x>1$).
  \item Compute $\dd{^{2}y}{x^{2}}$ implicitly for $x^{2}y=1$.
  \item Compute $\dd{y}{x}$ implicitly for $\cos(y)+x=y$.
  \item Find the equation of the tangent line to $y^{2}-x^{2}=3$ at
        $(1,2)$.
  \item Use logarithmic differentiation to find $y'$ for
        \[
          y=\frac{(2x+1)^{5}\sqrt{x^{2}+3}}{(x-4)^{3}}\qquad (x>4).
        \]
  \item Compute $\dd{y}{x}$ implicitly for $\mathrm{e}^{x+y}=xy$.
  \item Compute $\dd{^{2}y}{x^{2}}$ implicitly for $y^{3}=x^{2}$.
  \item Compute $\dd{y}{x}$ implicitly for $x^{2}+2xy+3y^{2}=12$.
  \item Find the equation of the tangent line to $\mathrm{e}^{y}+xy=\mathrm{e}+1$
        at $(1,1)$.
  \item Use logarithmic differentiation to find $y'$ for $y=(x^{2}+1)^{x}$
        ($x>0$).
  \item Compute $\dd{y}{x}$ implicitly for $\sqrt{x+y}=x-y$.
  \item Compute $\dd{^{2}y}{x^{2}}$ implicitly for $xy+y^{2}=1$.
  \item Find the equation of the tangent line to $x^{2}+y^{2}=4y$ at
        $(\sqrt{3},1)$.
  \item Compute $\dd{y}{x}$ implicitly for $\tan(xy)=x$.
  \item Use logarithmic differentiation to find $y'$ for $y=x^{1/x}$
        ($x>0$).
  \item Compute $\dd{y}{x}$ implicitly for $\ln(y)+xy=2$.
  \item Compute $\dd{y}{x}$ implicitly for $y^{3}+y=x^{2}$.
  \item Find the tangent line to the astroid $x^{2/3}+y^{2/3}=2$ at
        $(1,1)$.
  \item Compute $\dd{y}{x}$ implicitly for $\sin(x+y)=y$.
  \item Use logarithmic differentiation to find $y'$ for
        $y=\dfrac{\mathrm{e}^{2x}(x+1)^{3}}{\sqrt{x^{2}+5}}$ (assume $x>-1$).
  \item Compute $\dd{^{2}y}{x^{2}}$ implicitly for $x^{2}+y^{2}=25$.
  \item For an indifference curve given by $U(x,y)=x^{1/2}y^{1/2}=\bar U$
        (constant), compute $\dd{y}{x}$, the marginal rate of
        substitution.
  \item Find all points on the curve $x^{2}+y^{2}=2x+4y$ where the tangent
        line is horizontal.
  \item Compute $\dd{y}{x}$ implicitly for $\arctan(y/x)=x$.
  \item Use logarithmic differentiation to find $y'$ for $y=x^{x}$
        ($x>0$).
  \item Find the equation of the tangent line to $x^{3}+y^{3}=9xy$ at
        $(2,4)$ (point on the folium of Descartes).
  \item Compute $\dd{^{2}y}{x^{2}}$ implicitly for $y^{2}=4x$.
  \item Compute $\dd{y}{x}$ implicitly for $x\mathrm{e}^{y}+y=x$.
  \item Find all points on the ellipse $\dfrac{x^{2}}{9}+\dfrac{y^{2}}{4}=1$
        where the tangent line has slope $-1$.
  \item Use logarithmic differentiation to find $y'$ for
        $y=(\sin x)^{\cos x}$ on $(0,\pi/2)$.
  \item Compute $\dd{y}{x}$ implicitly for $x^{2}y-xy^{2}+x^{2}+y^{2}=0$.
  \item For the production isoquant $K^{1/3}L^{2/3}=\bar Q$ (constant),
        compute $\dd{K}{L}$, the marginal rate of technical
        substitution.
  \item Compute $\dd{^{2}y}{x^{2}}$ implicitly for $\sin(y)=x$.
  \item Find the tangent line to $xy=4$ at $x=2$.
  \item Use logarithmic differentiation to find $y'$ for
        $y=\dfrac{\sqrt{x+1}\sqrt[3]{x-2}}{(x+5)^{4}}$ (assume $x>2$).
  \item Compute $\dd{y}{x}$ implicitly for $\mathrm{e}^{y}\ln(x)+x^{2}=y$.
  \item Show that on the circle $x^{2}+y^{2}=r^{2}$ the tangent at
        $(x_{0},y_{0})$ is perpendicular to the radius. (I.e.\ derive the
        tangent slope and compare with the radius slope.)
  \item Use logarithmic differentiation to find $y'$ for $y=7^{x^{2}}$
        ($y>0$).
  \item Compute $\dd{^{2}y}{x^{2}}$ implicitly for $x^{2}+xy+y^{2}=3$.
  \item For the budget surface $px+qy=m$ (a line, $p,q,m>0$ constants),
        differentiate implicitly to recover $\dd{y}{x}=-p/q$.
  \item Find the tangent line to $x^{2}-xy+y^{2}=7$ at $(3,2)$.
  \item Compute $\dd{y}{x}$ implicitly for $\cos(xy)=x+y$.
  \item Compute $\dd{^{2}y}{x^{2}}$ implicitly for $x^{3}+y^{3}=6xy$ at the
        point $(3,3)$ (folium of Descartes). Express the answer as a
        number.
\end{problems}

\section{Solutions}

\begin{solutions}
  \item $2x+8yy'=0\Rightarrow y'=-\dfrac{x}{4y}$.

  \item Check $(2,2)$: $4+16=20$. From $2x+8yy'=0$: $y'=-x/(4y)$. At
        $(2,2)$: $y'=-2/8=-1/4$. Tangent: $y-2=-\tfrac14(x-2)$, i.e.\
        $x+4y=10$.

  \item $\ln y=\sin(x)\ln(x)\Rightarrow\dfrac{y'}{y}=\cos(x)\ln(x)+\dfrac{\sin x}{x}$.
        So
        \[
          y'=x^{\sin x}\left[\cos(x)\ln(x)+\frac{\sin x}{x}\right] .
        \]

  \item $2x+8yy'=0\Rightarrow y'=-\dfrac{x}{4y}$. Differentiate again:
        \[
          y''=-\frac{4y\cdot 1-x\cdot 4y'}{(4y)^{2}}=-\frac{y-xy'}{4y^{2}} .
        \]
        Substitute $y'=-x/(4y)$:
        $y-xy'=y+\dfrac{x^{2}}{4y}=\dfrac{4y^{2}+x^{2}}{4y}=\dfrac{4}{4y}
        =\dfrac{1}{y}$ (using $x^{2}+4y^{2}=4$). Hence
        $y''=-\dfrac{1}{4y^{3}}$.

  \item $3x^{2}-3y^{2}y'=3y+3xy'\Rightarrow -3(y^{2}+x)y'=3(y-x^{2})
        \Rightarrow y'=\dfrac{x^{2}-y}{y^{2}+x}$.

  \item $2xy^{3}+3x^{2}y^{2}y'=0\Rightarrow
        y'=-\dfrac{2xy^{3}}{3x^{2}y^{2}}=-\dfrac{2y}{3x}$.

  \item Differentiate $x^{2}y+y^{2}x=6$: $2xy+x^{2}y'+2yxy'+y^{2}=0
        \Rightarrow y'(x^{2}+2xy)=-(2xy+y^{2})
        \Rightarrow y'=-\dfrac{y(2x+y)}{x(x+2y)}$. At $(1,2)$:
        $y'=-\dfrac{2\cdot 4}{1\cdot 5}=-\dfrac{8}{5}$. Tangent:
        $y-2=-\tfrac85(x-1)$, i.e.\ $8x+5y=18$.

  \item $\ln y=x\ln(\ln x)\Rightarrow\dfrac{y'}{y}=\ln(\ln x)+x\cdot
        \dfrac{1/x}{\ln x}=\ln(\ln x)+\dfrac{1}{\ln x}$. So
        $y'=(\ln x)^{x}\left[\ln(\ln x)+\dfrac{1}{\ln x}\right]$.

  \item Explicit: $y=x^{-2}$, $y'=-2x^{-3}$, $y''=6x^{-4}$. Implicit
        check: $2xy+x^{2}y'=0\Rightarrow y'=-2y/x$. Differentiate:
        \[
          y''=-\frac{2y'x-2y}{x^{2}}=-\frac{2(-2y/x)x-2y}{x^{2}}
             =-\frac{-4y-2y}{x^{2}}=\frac{6y}{x^{2}}=\frac{6}{x^{4}} .
        \]

  \item $-\sin(y)y'+1=y'\Rightarrow y'(1+\sin y)=1\Rightarrow
        y'=\dfrac{1}{1+\sin y}$.

  \item $2yy'-2x=0\Rightarrow y'=x/y$. At $(1,2)$: $y'=1/2$. Tangent:
        $y-2=\tfrac12(x-1)$, i.e.\ $x-2y+3=0$.

  \item $\ln y=5\ln(2x+1)+\tfrac12\ln(x^{2}+3)-3\ln(x-4)$. Then
        \[
          \frac{y'}{y}=\frac{10}{2x+1}+\frac{x}{x^{2}+3}-\frac{3}{x-4} .
        \]
        Hence
        \[
          y'=\frac{(2x+1)^{5}\sqrt{x^{2}+3}}{(x-4)^{3}}
             \left[\frac{10}{2x+1}+\frac{x}{x^{2}+3}-\frac{3}{x-4}\right] .
        \]

  \item $\mathrm{e}^{x+y}(1+y')=y+xy'\Rightarrow
        y'(\mathrm{e}^{x+y}-x)=y-\mathrm{e}^{x+y}\Rightarrow
        y'=\dfrac{y-\mathrm{e}^{x+y}}{\mathrm{e}^{x+y}-x}$.

  \item $3y^{2}y'=2x\Rightarrow y'=\dfrac{2x}{3y^{2}}$. Differentiate
        again:
        \[
          y''=\frac{2\cdot 3y^{2}-2x\cdot 6yy'}{9y^{4}}
             =\frac{6y^{2}-12xy\cdot y'}{9y^{4}}
             =\frac{2}{3y^{2}}-\frac{4xy'}{3y^{3}} .
        \]
        Substitute $y'=2x/(3y^{2})$:
        \[
          y''=\frac{2}{3y^{2}}-\frac{4x}{3y^{3}}\cdot\frac{2x}{3y^{2}}
             =\frac{2}{3y^{2}}-\frac{8x^{2}}{9y^{5}}
             =\frac{6y^{3}-8x^{2}}{9y^{5}}
             =\frac{6x^{2}-8x^{2}}{9y^{5}}=-\frac{2x^{2}}{9y^{5}},
        \]
        using $y^{3}=x^{2}$. Equivalently, $y''=-\dfrac{2y^{3}}{9y^{5}}
        =-\dfrac{2}{9y^{2}}$.

  \item $2x+2y+2xy'+6yy'=0\Rightarrow y'(2x+6y)=-(2x+2y)\Rightarrow
        y'=-\dfrac{x+y}{x+3y}$.

  \item Differentiate $\mathrm{e}^{y}+xy=\mathrm{e}+1$:
        $\mathrm{e}^{y}y'+y+xy'=0\Rightarrow y'(\mathrm{e}^{y}+x)=-y
        \Rightarrow y'=-\dfrac{y}{\mathrm{e}^{y}+x}$. At $(1,1)$:
        $y'=-\dfrac{1}{\mathrm{e}+1}$. Tangent:
        $y-1=-\dfrac{1}{\mathrm{e}+1}(x-1)$.

  \item $\ln y=x\ln(x^{2}+1)\Rightarrow\dfrac{y'}{y}=\ln(x^{2}+1)
        +x\cdot\dfrac{2x}{x^{2}+1}=\ln(x^{2}+1)+\dfrac{2x^{2}}{x^{2}+1}$.
        So
        \[
          y'=(x^{2}+1)^{x}\left[\ln(x^{2}+1)+\frac{2x^{2}}{x^{2}+1}\right] .
        \]

  \item $\dfrac{1+y'}{2\sqrt{x+y}}=1-y'$. Multiply by $2\sqrt{x+y}$:
        $1+y'=2\sqrt{x+y}(1-y')\Rightarrow y'(1+2\sqrt{x+y})
        =2\sqrt{x+y}-1\Rightarrow y'=\dfrac{2\sqrt{x+y}-1}{1+2\sqrt{x+y}}$.

  \item $y+xy'+2yy'=0\Rightarrow y'(x+2y)=-y\Rightarrow
        y'=-\dfrac{y}{x+2y}$. Differentiate again using the quotient
        rule:
        \[
          y''=-\frac{y'(x+2y)-y(1+2y')}{(x+2y)^{2}}
             =-\frac{y'(x+2y)-y-2yy'}{(x+2y)^{2}}
             =-\frac{xy'-y}{(x+2y)^{2}} .
        \]
        Substitute $y'=-y/(x+2y)$:
        \[
          xy'-y=-\frac{xy}{x+2y}-y=-\frac{xy+y(x+2y)}{x+2y}
               =-\frac{2xy+2y^{2}}{x+2y}=-\frac{2y(x+y)}{x+2y} .
        \]
        Hence
        \[
          y''=\frac{2y(x+y)}{(x+2y)^{3}} .
        \]

  \item $2x+2yy'=4y'\Rightarrow y'(2y-4)=-2x\Rightarrow
        y'=\dfrac{x}{2-y}$. At $(\sqrt{3},1)$: $y'=\dfrac{\sqrt{3}}{1}
        =\sqrt{3}$. Tangent: $y-1=\sqrt{3}(x-\sqrt{3})$, i.e.\
        $y=\sqrt{3}x-2$.

  \item $\sec^{2}(xy)\cdot(y+xy')=1\Rightarrow y+xy'=\cos^{2}(xy)
        \Rightarrow y'=\dfrac{\cos^{2}(xy)-y}{x}$.

  \item $\ln y=\dfrac{\ln x}{x}\Rightarrow\dfrac{y'}{y}
        =\dfrac{(1/x)\cdot x-\ln x\cdot 1}{x^{2}}=\dfrac{1-\ln x}{x^{2}}$.
        So $y'=x^{1/x}\cdot\dfrac{1-\ln x}{x^{2}}$.

  \item $\dfrac{y'}{y}+y+xy'=0\Rightarrow y'\left(\dfrac1y+x\right)=-y
        \Rightarrow y'=-\dfrac{y^{2}}{1+xy}$.

  \item $3y^{2}y'+y'=2x\Rightarrow y'=\dfrac{2x}{3y^{2}+1}$.

  \item Differentiate $x^{2/3}+y^{2/3}=2$: $\tfrac23x^{-1/3}
        +\tfrac23y^{-1/3}y'=0\Rightarrow y'=-\dfrac{y^{1/3}}{x^{1/3}}
        =-\sqrt[3]{y/x}$. At $(1,1)$: $y'=-1$. Tangent: $y-1=-(x-1)$,
        i.e.\ $y=-x+2$.

  \item $\cos(x+y)(1+y')=y'\Rightarrow\cos(x+y)+\cos(x+y)y'=y'
        \Rightarrow\cos(x+y)=y'[1-\cos(x+y)]\Rightarrow
        y'=\dfrac{\cos(x+y)}{1-\cos(x+y)}$.

  \item $\ln y=2x+3\ln(x+1)-\tfrac12\ln(x^{2}+5)$. Differentiate:
        \[
          \frac{y'}{y}=2+\frac{3}{x+1}-\frac{x}{x^{2}+5} .
        \]
        Hence
        \[
          y'=\frac{\mathrm{e}^{2x}(x+1)^{3}}{\sqrt{x^{2}+5}}
             \left[2+\frac{3}{x+1}-\frac{x}{x^{2}+5}\right] .
        \]

  \item $2x+2yy'=0\Rightarrow y'=-x/y$. Differentiate again:
        \[
          y''=-\frac{y\cdot 1-x\cdot y'}{y^{2}}
             =-\frac{y-x(-x/y)}{y^{2}}
             =-\frac{y+x^{2}/y}{y^{2}}
             =-\frac{y^{2}+x^{2}}{y^{3}}=-\frac{25}{y^{3}},
        \]
        using $x^{2}+y^{2}=25$.

  \item Differentiate $x^{1/2}y^{1/2}=\bar U$:
        $\tfrac12x^{-1/2}y^{1/2}+x^{1/2}\cdot\tfrac12y^{-1/2}y'=0$.
        Multiply by $2x^{1/2}y^{1/2}$: $y+xy'=0\Rightarrow y'=-y/x$. So
        $\dd{y}{x}=-\dfrac{y}{x}$ (negative slope along the indifference
        curve, standard MRS).

  \item $2x+2yy'=2+4y'\Rightarrow y'(2y-4)=2-2x\Rightarrow
        y'=\dfrac{1-x}{y-2}$. Horizontal tangent: $y'=0$ requires $x=1$
        (and $y\neq 2$). Plug into the curve: $1+y^{2}=2+4y\Rightarrow
        y^{2}-4y-1=0\Rightarrow y=2\pm\sqrt{5}$. So horizontal tangents at
        $(1,2+\sqrt{5})$ and $(1,2-\sqrt{5})$.

  \item Differentiate $\arctan(y/x)=x$:
        $\dfrac{1}{1+(y/x)^{2}}\cdot\dfrac{y'x-y}{x^{2}}=1$, i.e.\
        $\dfrac{y'x-y}{x^{2}+y^{2}}=1$. Hence $y'x-y=x^{2}+y^{2}$, so
        $y'=\dfrac{x^{2}+y^{2}+y}{x}$.

  \item $\ln y=x\ln x\Rightarrow\dfrac{y'}{y}=\ln x+1\Rightarrow
        y'=x^{x}(\ln x+1)$.

  \item Check $(2,4)$ on $x^{3}+y^{3}=9xy$: $8+64=72$, $9\cdot 8=72$.
        Good. Differentiate: $3x^{2}+3y^{2}y'=9y+9xy'\Rightarrow
        y'(3y^{2}-9x)=9y-3x^{2}\Rightarrow
        y'=\dfrac{9y-3x^{2}}{3y^{2}-9x}=\dfrac{3y-x^{2}}{y^{2}-3x}$.
        At $(2,4)$: $y'=\dfrac{12-4}{16-6}=\dfrac{8}{10}=\dfrac{4}{5}$.
        Tangent: $y-4=\tfrac45(x-2)$, i.e.\ $4x-5y+12=0$.

  \item $2yy'=4\Rightarrow y'=2/y$. Differentiate again:
        $y''=-\dfrac{2y'}{y^{2}}=-\dfrac{2(2/y)}{y^{2}}=-\dfrac{4}{y^{3}}$.

  \item $\mathrm{e}^{y}+x\mathrm{e}^{y}y'+y'=1\Rightarrow
        y'(x\mathrm{e}^{y}+1)=1-\mathrm{e}^{y}\Rightarrow
        y'=\dfrac{1-\mathrm{e}^{y}}{1+x\mathrm{e}^{y}}$.

  \item $\dfrac{2x}{9}+\dfrac{2y}{4}y'=0\Rightarrow y'=-\dfrac{4x}{9y}$.
        Set $-\dfrac{4x}{9y}=-1\Rightarrow 4x=9y\Rightarrow y=\dfrac{4x}{9}$.
        Plug into the ellipse: $\dfrac{x^{2}}{9}+\dfrac{(4x/9)^{2}}{4}=1
        \Rightarrow\dfrac{x^{2}}{9}+\dfrac{16x^{2}}{4\cdot 81}=1
        \Rightarrow\dfrac{x^{2}}{9}+\dfrac{4x^{2}}{81}=1
        \Rightarrow\dfrac{9x^{2}+4x^{2}}{81}=1\Rightarrow
        13x^{2}=81\Rightarrow x=\pm 9/\sqrt{13}$. Corresponding
        $y=\pm 4/\sqrt{13}$ (same sign as $x$). Points:
        $\left(\dfrac{9}{\sqrt{13}},\dfrac{4}{\sqrt{13}}\right)$ and
        $\left(-\dfrac{9}{\sqrt{13}},-\dfrac{4}{\sqrt{13}}\right)$.

  \item $\ln y=\cos(x)\ln(\sin x)\Rightarrow\dfrac{y'}{y}
        =-\sin(x)\ln(\sin x)+\cos(x)\cdot\dfrac{\cos x}{\sin x}
        =-\sin(x)\ln(\sin x)+\dfrac{\cos^{2}x}{\sin x}$. So
        \[
          y'=(\sin x)^{\cos x}\left[-\sin(x)\ln(\sin x)
             +\frac{\cos^{2}x}{\sin x}\right] .
        \]

  \item $2xy+x^{2}y'-y^{2}-x\cdot 2yy'+2x+2yy'=0$. Collect $y'$:
        $y'(x^{2}-2xy+2y)=y^{2}-2xy-2x$. So
        $y'=\dfrac{y^{2}-2xy-2x}{x^{2}-2xy+2y}$.

  \item $\tfrac13K^{-2/3}L^{2/3}\dd{K}{L}\cdot L^{2/3}\ldots$
        carefully: differentiate $K^{1/3}L^{2/3}=\bar Q$ with respect to
        $L$, treating $K=K(L)$:
        \[
          \tfrac13K^{-2/3}K'(L)\cdot L^{2/3}+K^{1/3}\cdot\tfrac23L^{-1/3}=0 .
        \]
        Multiply by $3K^{2/3}L^{1/3}$: $LK'(L)+2K=0\Rightarrow
        K'(L)=-\dfrac{2K}{L}$. So $\dd{K}{L}=-\dfrac{2K}{L}$ (the MRTS).

  \item $\cos(y)y'=1\Rightarrow y'=\dfrac{1}{\cos y}=\sec(y)$.
        Differentiate again: $y''=\sec(y)\tan(y)\cdot y'
        =\sec^{2}(y)\tan(y)$.

  \item At $x=2$, $y=2$. Differentiate $xy=4$: $y+xy'=0\Rightarrow
        y'=-y/x$. At $(2,2)$: $y'=-1$. Tangent: $y-2=-(x-2)$, i.e.\
        $y=-x+4$.

  \item $\ln y=\tfrac12\ln(x+1)+\tfrac13\ln(x-2)-4\ln(x+5)$.
        Differentiate:
        \[
          \frac{y'}{y}=\frac{1}{2(x+1)}+\frac{1}{3(x-2)}-\frac{4}{x+5} .
        \]
        Hence
        \[
          y'=\frac{\sqrt{x+1}\sqrt[3]{x-2}}{(x+5)^{4}}
             \left[\frac{1}{2(x+1)}+\frac{1}{3(x-2)}-\frac{4}{x+5}\right] .
        \]

  \item $\mathrm{e}^{y}y'\ln(x)+\dfrac{\mathrm{e}^{y}}{x}+2x=y'\Rightarrow
        y'[\mathrm{e}^{y}\ln(x)-1]=-\dfrac{\mathrm{e}^{y}}{x}-2x\Rightarrow
        y'=\dfrac{\mathrm{e}^{y}/x+2x}{1-\mathrm{e}^{y}\ln(x)}$.

  \item Differentiate $x^{2}+y^{2}=r^{2}$: $2x+2yy'=0\Rightarrow
        y'=-x/y$. The radius from origin to $(x_{0},y_{0})$ has slope
        $y_{0}/x_{0}$. Product of slopes:
        $(-x_{0}/y_{0})\cdot(y_{0}/x_{0})=-1$, confirming
        perpendicularity.

  \item $\ln y=x^{2}\ln(7)\Rightarrow\dfrac{y'}{y}=2x\ln(7)\Rightarrow
        y'=2x\ln(7)\cdot 7^{x^{2}}$.

  \item Differentiate $x^{2}+xy+y^{2}=3$: $2x+y+xy'+2yy'=0\Rightarrow
        y'(x+2y)=-(2x+y)\Rightarrow y'=-\dfrac{2x+y}{x+2y}$.
        Differentiate again using the quotient rule:
        \[
          y''=-\frac{(2+y')(x+2y)-(2x+y)(1+2y')}{(x+2y)^{2}} .
        \]
        Expand the numerator: $(2+y')(x+2y)=2x+4y+xy'+2yy'$, and
        $(2x+y)(1+2y')=2x+y+4xy'+2yy'$. Subtract:
        $N:=(2+y')(x+2y)-(2x+y)(1+2y')=(2x+4y+xy'+2yy')
        -(2x+y+4xy'+2yy')=3y-3xy'$. So
        $y''=-\dfrac{3y-3xy'}{(x+2y)^{2}}=\dfrac{3(xy'-y)}{(x+2y)^{2}}$.
        Substitute $y'=-\dfrac{2x+y}{x+2y}$:
        \[
          xy'-y=-\frac{x(2x+y)}{x+2y}-y
               =-\frac{x(2x+y)+y(x+2y)}{x+2y}
        \]
        \[
          \phantom{xy'-y}=-\frac{2x^{2}+xy+xy+2y^{2}}{x+2y}
               =-\frac{2(x^{2}+xy+y^{2})}{x+2y},
        \]
        using $x^{2}+xy+y^{2}=3$. Hence
        \[
          y''=\frac{3}{(x+2y)^{2}}\cdot\left(-\frac{6}{x+2y}\right)
             =-\frac{18}{(x+2y)^{3}} .
        \]

  \item Differentiate $px+qy=m$ w.r.t.\ $x$: $p+qy'=0\Rightarrow
        y'=-p/q$, as required (the budget line has constant slope).

  \item Differentiate $x^{2}-xy+y^{2}=7$: $2x-y-xy'+2yy'=0\Rightarrow
        y'(2y-x)=y-2x\Rightarrow y'=\dfrac{y-2x}{2y-x}$. At $(3,2)$:
        $y'=\dfrac{2-6}{4-3}=-4$. Tangent: $y-2=-4(x-3)$, i.e.\
        $y=-4x+14$.

  \item Differentiate $\cos(xy)=x+y$: $-\sin(xy)\cdot(y+xy')=1+y'
        \Rightarrow -y\sin(xy)-x\sin(xy)y'=1+y'\Rightarrow
        -y'(x\sin(xy)+1)=1+y\sin(xy)\Rightarrow
        y'=-\dfrac{1+y\sin(xy)}{1+x\sin(xy)}$.

  \item First verify $(3,3)$: $27+27=54$ and $6\cdot 9=54$. Good.
        Differentiate $x^{3}+y^{3}=6xy$: $3x^{2}+3y^{2}y'=6y+6xy'
        \Rightarrow y'(3y^{2}-6x)=6y-3x^{2}\Rightarrow
        y'=\dfrac{6y-3x^{2}}{3y^{2}-6x}=\dfrac{2y-x^{2}}{y^{2}-2x}$.
        At $(3,3)$: $y'=\dfrac{6-9}{9-6}=-1$. For $y''$, differentiate
        $y'=\dfrac{2y-x^{2}}{y^{2}-2x}$:
        \[
          y''=\frac{(2y'-2x)(y^{2}-2x)-(2y-x^{2})(2yy'-2)}{(y^{2}-2x)^{2}} .
        \]
        At $(3,3)$ with $y'=-1$: numerator pieces:
        \[
          (2y'-2x)(y^{2}-2x)=(-2-6)(9-6)=(-8)(3)=-24,
        \]
        \[
          (2y-x^{2})(2yy'-2)=(6-9)(6\cdot(-1)-2)=(-3)(-8)=24 .
        \]
        So numerator $=-24-24=-48$ and denominator $=(9-6)^{2}=9$.
        Hence $y''(3)=-\dfrac{48}{9}=-\dfrac{16}{3}$.
\end{solutions}

\section{Additional Problems}

\begin{problems}
  \item Compute $\dd{y}{x}$ implicitly for $x^{2}+9y^{2}=36$.
  \item Find the equation of the tangent line to $x^{2}+y^{2}=25$ at
        $(3,4)$.
  \item Use logarithmic differentiation to find $y'$ for $y=x^{\cos x}$
        ($x>0$).
  \item Compute $\dd{y}{x}$ implicitly for $x^{2}y+xy^{2}=6$.
  \item Compute $\dd{^{2}y}{x^{2}}$ implicitly for $x^{2}-y^{2}=4$.
  \item Find the equation of the tangent line to $x^{3}+xy+y^{3}=11$ at
        $(1,2)$.
  \item Use logarithmic differentiation to find $y'$ for
        \[
          y=\frac{(x+1)^{2}\sqrt{x-3}}{(2x+5)^{4}}\qquad (x>3).
        \]
  \item Compute $\dd{y}{x}$ implicitly for $\mathrm{e}^{xy}=x+y$.
  \item For an indifference curve $U(x,y)=x^{1/3}y^{2/3}=\bar U$
        (constant), compute $\dd{y}{x}$, the marginal rate of
        substitution.
  \item Compute $\dd{y}{x}$ implicitly for $\sin(x+y)=xy$.
  \item Find the equation of the tangent line to $x^{2}-2xy+3y^{2}=9$ at
        $(1,2)$.
  \item Compute $\dd{^{2}y}{x^{2}}$ implicitly for $x^{2}+y^{2}=9$.
  \item Find all points on the curve $x^{2}+y^{2}=4x-2y$ where the tangent
        line is vertical.
  \item For the production isoquant $K^{2/5}L^{3/5}=\bar Q$ (constant),
        compute $\dd{K}{L}$, the marginal rate of technical
        substitution.
  \item Use logarithmic differentiation to find $y'$ for
        $y=(x^{2}+1)^{\sin x}$ ($x\in\R$).

  \item The curve $k$ below is the set of points satisfying
        \[
          x^{3}-3x^{2}+y^{2}-2y=-1 .
        \]
        Verify that $(0,1)$, $(3,1)$, $(-1,3)$ and $(-1,-1)$ all lie on
        $k$.

\begin{gallery}
\begin{tikzpicture}
  \begin{axis}[geoaxis, xmin=-1.6, xmax=3.6, ymin=-1.9, ymax=3.9,
      xtick={-1,1,2,3}, ytick={-1,1,2,3}, title={the curve $k$}]
    \addplot[domain=-1.2:3, samples=241]{1+abs(x)*sqrt(3-x)};
    \addplot[domain=-1.2:3, samples=241]{1-abs(x)*sqrt(3-x)};
  \end{axis}
\end{tikzpicture}\hfill
\begin{tikzpicture}
  \begin{axis}[geoaxis, xmin=-1.4, xmax=1.4, ymin=-0.9, ymax=0.9,
      xtick={-1,1}, ytick={-0.5,0.5}, title={the curve $\ell$}]
    \addplot[domain=-1:1, samples=241]{x*sqrt(1-x^2)};
    \addplot[domain=-1:1, samples=241]{-x*sqrt(1-x^2)};
  \end{axis}
\end{tikzpicture}
\end{gallery}

  \item Name every point of $k$ at which the hypotheses of the implicit
        function theorem fail for writing $y$ as a function of $x$, and say
        which graphical feature identifies each one.
  \item Confirm the previous answer by computation: solve $F_{y}=0$
        together with $F=0$.
  \item Compute the slope of the tangent to $k$ at $(-1,3)$.
  \item Compute the slope of the tangent to $k$ at $(-1,-1)$, and explain
        why it is the negative of the previous answer.
  \item Find every point of $k$ at which the tangent is horizontal.
  \item Near $(3,1)$ no function $y(x)$ describes $k$. Show that
        $x=x(y)$ does exist there, and compute $\dd{x}{y}$ at that point.
  \item The curve $\ell$ above is $y^{2}=x^{2}\bigl(1-x^{2}\bigr)$. Write it
        as $F(x,y)=0$ and give $F_{x}$ and $F_{y}$.
  \item Name every point of $\ell$ at which the implicit function theorem
        fails, and the graphical feature of each.
  \item Compute the slope of the tangent to $\ell$ at
        $\bigl(\tfrac12,\tfrac{\sqrt3}{4}\bigr)$.
  \item Compute the slope of the tangent to $\ell$ at
        $\bigl(\tfrac12,-\tfrac{\sqrt3}{4}\bigr)$.
  \item Find the four points of $\ell$ at which the tangent is horizontal.
  \item At the origin two branches of $\ell$ cross. Separate them by solving
        for $y$, and compute the slope of each branch there.
  \item Explain in one or two sentences why $-F_{x}/F_{y}$ cannot deliver
        those two slopes, and what has to be done instead.
\end{problems}

\section{Additional Solutions}

\begin{solutions}
  \item $2x+18yy'=0\Rightarrow y'=-\dfrac{x}{9y}$.

  \item Check $(3,4)$: $9+16=25$. From $2x+2yy'=0$: $y'=-x/y$. At
        $(3,4)$: $y'=-3/4$. Tangent: $y-4=-\tfrac34(x-3)$, i.e.\
        $3x+4y=25$.

  \item $\ln y=\cos(x)\ln(x)\Rightarrow\dfrac{y'}{y}
        =-\sin(x)\ln(x)+\dfrac{\cos x}{x}$. So
        \[
          y'=x^{\cos x}\left[-\sin(x)\ln(x)+\frac{\cos x}{x}\right] .
        \]

  \item Differentiate $x^{2}y+xy^{2}=6$: $2xy+x^{2}y'+y^{2}+2xyy'=0
        \Rightarrow y'(x^{2}+2xy)=-(2xy+y^{2})
        \Rightarrow y'=-\dfrac{y(2x+y)}{x(x+2y)}$.

  \item $2x-2yy'=0\Rightarrow y'=x/y$. Differentiate again:
        \[
          y''=\frac{y-xy'}{y^{2}}=\frac{y-x(x/y)}{y^{2}}
             =\frac{y^{2}-x^{2}}{y^{3}}=\frac{-4}{y^{3}},
        \]
        using $x^{2}-y^{2}=4\Rightarrow y^{2}-x^{2}=-4$.

  \item Check $(1,2)$: $1+2+8=11$. Differentiate:
        $3x^{2}+y+xy'+3y^{2}y'=0\Rightarrow y'(x+3y^{2})=-(3x^{2}+y)
        \Rightarrow y'=-\dfrac{3x^{2}+y}{x+3y^{2}}$. At $(1,2)$:
        $y'=-\dfrac{3+2}{1+12}=-\dfrac{5}{13}$. Tangent:
        $y-2=-\tfrac{5}{13}(x-1)$, i.e.\ $5x+13y=31$.

  \item $\ln y=2\ln(x+1)+\tfrac12\ln(x-3)-4\ln(2x+5)$. Then
        \[
          \frac{y'}{y}=\frac{2}{x+1}+\frac{1}{2(x-3)}-\frac{8}{2x+5} .
        \]
        Hence
        \[
          y'=\frac{(x+1)^{2}\sqrt{x-3}}{(2x+5)^{4}}
             \left[\frac{2}{x+1}+\frac{1}{2(x-3)}-\frac{8}{2x+5}\right] .
        \]

  \item Differentiate $\mathrm{e}^{xy}=x+y$:
        $\mathrm{e}^{xy}(y+xy')=1+y'\Rightarrow
        y'(x\mathrm{e}^{xy}-1)=1-y\mathrm{e}^{xy}\Rightarrow
        y'=\dfrac{1-y\mathrm{e}^{xy}}{x\mathrm{e}^{xy}-1}$.

  \item Differentiate $x^{1/3}y^{2/3}=\bar U$:
        $\tfrac13x^{-2/3}y^{2/3}+x^{1/3}\cdot\tfrac23y^{-1/3}y'=0$.
        Multiply by $3x^{2/3}y^{1/3}$: $y+2xy'=0\Rightarrow y'=-y/(2x)$.
        So $\dd{y}{x}=-\dfrac{y}{2x}$ (the MRS).

  \item Differentiate $\sin(x+y)=xy$: $\cos(x+y)(1+y')=y+xy'
        \Rightarrow y'(\cos(x+y)-x)=y-\cos(x+y)\Rightarrow
        y'=\dfrac{y-\cos(x+y)}{\cos(x+y)-x}$.

  \item Check $(1,2)$: $1-4+12=9$. Good. Differentiate:
        $2x-2y-2xy'+6yy'=0\Rightarrow y'(6y-2x)=2y-2x\Rightarrow
        y'=\dfrac{y-x}{3y-x}$. At $(1,2)$: $y'=\dfrac{1}{5}$. Tangent:
        $y-2=\tfrac15(x-1)$, i.e.\ $x-5y+9=0$.

  \item $2x+2yy'=0\Rightarrow y'=-x/y$. Differentiate again:
        \[
          y''=-\frac{y-xy'}{y^{2}}=-\frac{y-x(-x/y)}{y^{2}}
             =-\frac{y+x^{2}/y}{y^{2}}
             =-\frac{y^{2}+x^{2}}{y^{3}}=-\frac{9}{y^{3}},
        \]
        using $x^{2}+y^{2}=9$.

  \item $2x+2yy'=4-2y'\Rightarrow y'(2y+2)=4-2x\Rightarrow
        y'=\dfrac{2-x}{y+1}$. Vertical tangent: $y'$ undefined when
        $y+1=0\Rightarrow y=-1$ (and $x\neq 2$). Plug into the curve:
        $x^{2}+1=4x+2\Rightarrow x^{2}-4x-1=0\Rightarrow
        x=2\pm\sqrt{5}$. Vertical tangents at $(2+\sqrt{5},-1)$ and
        $(2-\sqrt{5},-1)$.

  \item Differentiate $K^{2/5}L^{3/5}=\bar Q$ w.r.t.\ $L$, treating
        $K=K(L)$:
        \[
          \tfrac25K^{-3/5}K'(L)\cdot L^{3/5}+K^{2/5}\cdot\tfrac35L^{-2/5}=0 .
        \]
        Multiply by $5K^{3/5}L^{2/5}$: $2LK'(L)+3K=0\Rightarrow
        K'(L)=-\dfrac{3K}{2L}$. So $\dd{K}{L}=-\dfrac{3K}{2L}$ (the
        MRTS).

  \item $\ln y=\sin(x)\ln(x^{2}+1)\Rightarrow\dfrac{y'}{y}
        =\cos(x)\ln(x^{2}+1)+\sin(x)\cdot\dfrac{2x}{x^{2}+1}$. So
        \[
          y'=(x^{2}+1)^{\sin x}\left[\cos(x)\ln(x^{2}+1)
             +\frac{2x\sin x}{x^{2}+1}\right] .
        \]

  \item With $F(x,y)=x^{3}-3x^{2}+y^{2}-2y+1$:
        $F(0,1)=0-0+1-2+1=0$; \;$F(3,1)=27-27+1-2+1=0$;
        \;$F(-1,3)=-1-3+9-6+1=0$; \;$F(-1,-1)=-1-3+1+2+1=0$.
        All four are on $k$.
  \item Graphically: at $(0,1)$ the curve \textbf{crosses itself}, and at
        $(3,1)$ it has a \textbf{vertical tangent}. At every other point of
        $k$ a local $y(x)$ exists.
  \item $F_{y}=2y-2=0\iff y=1$. Substituting $y=1$ into $F=0$ gives
        $x^{3}-3x^{2}+1-2+1=x^{2}(x-3)=0$, so $x=0$ (a double root) or
        $x=3$: the points $(0,1)$ and $(3,1)$. At $(3,1)$,
        $F_{x}=3x^{2}-6x=9\neq0$, so this is the vertical-tangent case; at
        $(0,1)$, $F_{x}=0$ as well, the singular case.
  \item $F_{x}=3x^{2}-6x$, $F_{y}=2y-2$, so
        $y'=-\dfrac{3x^{2}-6x}{2y-2}$. At $(-1,3)$:
        $y'=-\dfrac{3+6}{6-2}=-\dfrac94$.
  \item At $(-1,-1)$: $y'=-\dfrac{3+6}{-2-2}=\dfrac94$. The equation
        contains $y$ only through $(y-1)^{2}$, so $k$ is symmetric about the
        line $y=1$; reflecting a curve in a horizontal line negates every
        slope, and $(-1,3)$ and $(-1,-1)$ are mirror images.
  \item Horizontal tangent needs $F_{x}=0$ and $F_{y}\neq0$:
        $3x^{2}-6x=3x(x-2)=0$ gives $x=0$ or $x=2$. The value $x=0$ forces
        $y=1$, where $F_{y}=0$ too --- that is the singular point, not a
        horizontal tangent. For $x=2$: $(y-1)^{2}=x^{2}(3-x)=4$, so $y=3$ or
        $y=-1$. The tangent is horizontal at $(2,3)$ and $(2,-1)$.
  \item At $(3,1)$ we have $F_{x}=9\neq0$, so the theorem applies with the
        roles of the variables exchanged and gives a function $x=x(y)$.
        Differentiating $F\bigl(x(y),y\bigr)=0$ with respect to $y$:
        $F_{x}\,x'(y)+F_{y}=0$, hence
        $\dd{x}{y}=-\dfrac{F_{y}}{F_{x}}=-\dfrac{0}{9}=0$ --- the curve is
        momentarily vertical, which is exactly a horizontal tangent for
        $x(y)$.
  \item $F(x,y)=y^{2}-x^{2}+x^{4}$, with $F_{x}=-2x+4x^{3}=2x(2x^{2}-1)$
        and $F_{y}=2y$.
  \item $F_{y}=2y=0\iff y=0$, and then $x^{2}(1-x^{2})=0$ gives
        $x=0,\pm1$. At $(\pm1,0)$: $F_{x}=\pm2\neq0$, \textbf{vertical
        tangents} --- the two tips of the figure eight. At $(0,0)$:
        $F_{x}=0$ too, the \textbf{self-intersection} in the middle.
  \item $y'=-\dfrac{2x(2x^{2}-1)}{2y}$. At
        $\bigl(\tfrac12,\tfrac{\sqrt3}{4}\bigr)$ the numerator is
        $2\cdot\tfrac12\bigl(\tfrac12-1\bigr)=-\tfrac12$ and the denominator
        is $\tfrac{\sqrt3}{2}$, so
        $y'=\dfrac{1/2}{\sqrt3/2}=\dfrac{1}{\sqrt3}=\dfrac{\sqrt3}{3}$.
  \item Only the sign of $y$ changes, so
        $y'=-\dfrac{\sqrt3}{3}$.
  \item $F_{x}=2x(2x^{2}-1)=0$ with $F_{y}\neq0$ gives $x=\pm
        \tfrac{1}{\sqrt2}$ (the root $x=0$ is the singular point). Then
        $y^{2}=x^{2}(1-x^{2})=\tfrac12\cdot\tfrac12=\tfrac14$, so
        $y=\pm\tfrac12$: the four points
        $\bigl(\pm\tfrac{1}{\sqrt2},\pm\tfrac12\bigr)$.
  \item Solving gives $y=x\sqrt{1-x^{2}}$ and $y=-x\sqrt{1-x^{2}}$. For the
        first,
        $y'=\sqrt{1-x^{2}}-\dfrac{x^{2}}{\sqrt{1-x^{2}}}$, which at $x=0$
        equals $1$; the second branch gives $-1$. The two branches cross the
        origin at $+45^{\circ}$ and $-45^{\circ}$.
  \item At the origin both partials vanish, so $-F_{x}/F_{y}$ is $0/0$: it
        carries no information, and a single number could not describe two
        different slopes anyway. The branches have to be separated first ---
        here by solving for $y$ --- and each differentiated on its own.
\end{solutions}
